\documentclass[ignorenonframetext,]{beamer}

%Set up notes
\setbeamertemplate{note page}[plain]
\setbeameroption{hide notes}

\usepackage{amssymb,amsmath}
\usepackage{ifxetex,ifluatex}
\usepackage{fixltx2e} % provides \textsubscript
\usepackage{tikz}
\usepackage{tikz-qtree}
\usepackage{pdfpages} %include beamer pdf slides from other presentations \usepackage{color}
\usepackage{hyperref}
\ifxetex \usepackage{fontspec,xltxtra,xunicode}
  \defaultfontfeatures{Mapping=tex-text,Scale=MatchLowercase}
\else
  \ifluatex
    \usepackage{fontspec}
    \defaultfontfeatures{Mapping=tex-text,Scale=MatchLowercase}
  \else
    \usepackage[utf8]{inputenc}
  \fi
\fi

% Comment these out if you don't want a slide with just the
% part/section/subsection/subsubsection title:
\AtBeginPart{
  \let\insertpartnumber\relax
  \let\partname\relax
  \frame{\partpage}
}
\AtBeginSection{
  \let\insertsectionnumber\relax
  \let\sectionname\relax
  \frame{\sectionpage}
}
\AtBeginSubsection{
  \let\insertsubsectionnumber\relax
  \let\subsectionname\relax
  \frame{\subsectionpage}
}

\setlength{\parindent}{0pt}
\setlength{\parskip}{6pt plus 2pt minus 1pt}
\setlength{\emergencystretch}{3em}  % prevent overfull lines
\setcounter{secnumdepth}{0}

\title{Bigger is better: Increasing Returns to Scale and Trade}
\author{Instructor: David Jinkins}
\date{Date: Sept. 18, 2014}

\begin{document}

\frame{\titlepage}

\begin{frame}
\begin{itemize}
\itemsep1pt\parskip0pt\parsep0pt
\item
  Last Time: The Standard Model 
  \begin{itemize}
        \item Model
        \begin{itemize}
            \item Focus on PPF and Demand
            \item Leave production details unspecified
            \item More careful about Demand
        \end{itemize}
        \item Growth
        \begin{itemize}
            \item Export-biased growth is bad
            \item Import-biased growth is good
            \item Theoretically possible to be hurt by growth
        \end{itemize}
        \item Tariffs and subsidies
        \begin{itemize}
            \item Tariffs on import good improve terms of trade
            \item Export subsidy worsens terms of trade
            \item Effect on welfare qualified
        \end{itemize}
        \item International lending
    \end{itemize}
\end{itemize}

\end{frame}

\begin{frame}
\begin{itemize}
\itemsep1pt\parskip0pt\parsep0pt
\item
  This time: Increasing Returns to Scale
  \begin{itemize}
        \item Krugman: External Economies 
        \begin{itemize}
            \item Larger industries have lower cost
            \item Drives industries to concentrate
            \item A reason for trade
        \end{itemize}
        \item Implications of external economies 
        \begin{itemize}
            \item Historical Accident
            \item Money on the table
            \item Infant industries
        \end{itemize}
        \item Melitz and Krugman: Internal Economies
        \begin{itemize}
            \item Larger firms have lower cost 
            \item Each firm a different product or variety
            \item Consumers like a mix
            \item Reason for trade: more (and cheaper) varieties
        \end{itemize}
        \item Applications
        \begin{itemize}
            \item Dumping
            \item Outsourcing
            \item FDI
        \end{itemize}
    \end{itemize}
\end{itemize}

\end{frame}

\begin{frame}
    \begin{itemize}
        \item But first a review!
    \end{itemize}
\end{frame}


\frame[plain]{\includegraphics[page=7,width=\textwidth]{introduction_gravity.pdf}}
\frame[plain]{\includegraphics[page=8,width=\textwidth]{introduction_gravity.pdf}}
\frame[plain]{\includegraphics[page=25,width=\textwidth]{standard_model.pdf}}
\frame[plain]{\includegraphics[page=30,width=\textwidth]{standard_model.pdf}}
\frame[plain]{\includegraphics[page=36,width=\textwidth]{standard_model.pdf}}
\frame[plain]{\includegraphics[page=37,width=\textwidth]{standard_model.pdf}}
\frame[plain]{\includegraphics[page=38,width=\textwidth]{standard_model.pdf}}
\frame[plain]{\includegraphics[page=40,width=\textwidth]{standard_model.pdf}}
\frame[plain]{\includegraphics[page=43,width=\textwidth]{standard_model.pdf}}
\frame[plain]{\includegraphics[page=57,width=\textwidth]{standard_model.pdf}}
\frame[plain]{\includegraphics[page=58,width=\textwidth]{standard_model.pdf}}
\frame[plain]{\includegraphics[page=67,width=\textwidth]{standard_model.pdf}}
\frame[plain]{\includegraphics[page=68,width=\textwidth]{standard_model.pdf}}
\frame[plain]{\includegraphics[page=69,width=\textwidth]{standard_model.pdf}}
\frame[plain]{\includegraphics[page=74,width=\textwidth]{standard_model.pdf}}
\frame[plain]{\includegraphics[page=76,width=\textwidth]{standard_model.pdf}}
\frame[plain]{\includegraphics[page=77,width=\textwidth]{standard_model.pdf}}
\frame[plain]{\includegraphics[page=78,width=\textwidth]{standard_model.pdf}}

\begin{frame}
    \begin{itemize}
        \item End review!
    \end{itemize}
\end{frame}

\begin{frame}{Chapter 7: External Economies}

    \begin{itemize}
        \item A trade model quite different from those earlier 
            \begin{itemize}
            \item Technology depends upon scale and experience
            \item Best to concentrate production in one location
            \item Drives countries to specialize, \emph{even if ex-ante identical!}
        \end{itemize}
        \item Idea: Larger industries have lower costs
        \begin{itemize}
            \item knowledge spillovers due to informal interaction
            \item Specialized Labor market pooling
            \item Specialized suppliers
        \end{itemize}
        \item Room for helpful government policy 
        \begin{itemize}
            \item Historical accident: Industries can be in the 'wrong' country 
            \item There can be equilibrium losses from trade!
            \item 'Infant' Industries may need early protection from competition
        \end{itemize}
    \end{itemize}
            
\end{frame}

\begin{frame}{Returns to Scale}
    \begin{itemize}
        \item IKEA furniture
        \begin{itemize}
            \item One person: All day to assemble a flipping dresser
            \item Two people: Two hours for a dresser
            \item Ten people: Can assemble 5 dressers in an hour 
            \item Factory of fifty people: Can assemble 400 dressers in eight hours
        \end{itemize}
        \item How many hours per dresser?
    \end{itemize}

\end{frame}

\begin{frame}{External vs Internal Increasing Returns}
    \begin{itemize}
        \item Firms or locations?
        \begin{itemize}
            \item If increasing returns happen within an industry and firm, what would the economy look like?
            \item If increasing returns happen within an industry and location, what would the economy look like? 
        \end{itemize}
        \item First: increasing returns within industry and location
        \begin{itemize}
            \item Firms are very small, no affect on aggregate production 
            \item The more the industry produces, the lower are average costs
            \item Good model for Silicon Valley startups
        \end{itemize}
    \end{itemize}
\end{frame}
\begin{frame}{Decreasing marginal product of labor?}

    \begin{itemize}
        \item Specific factors, HO -- increase labor, fix capital, decreasing returns to labor
        \item Now we want increasing returns -- what is the argument?
        \begin{itemize}
            \item Specialized suppliers
            \item Labor market pooling
            \item Knowledge spillovers
        \end{itemize}
    \end{itemize}

\end{frame}

\begin{frame}{Specialized suppliers}

    \begin{itemize}
        \item Production can require very specific inputs 
        \item Maybe: 
        \begin{itemize}
            \item Specialized tool for cleaning button making machines
            \item High quality polymers for production of CPU's
            \item Chemicals used in pharmaceutical research
        \end{itemize}
        \item If industry is clusters, suppliers likely to cluster as well
        \item Costly to open new firm in completely different location
    \end{itemize}

\end{frame}

\begin{frame}{Labor market pooling}

    \begin{itemize}
        \item We have been treating labor as costlessly mobile (equal wages)
        \item Now let's make it totally impossible for labor to move between locations
        \begin{itemize}
            \item Suppose we have two universities, one in Aarhus and one in Copenhagen
            \item Suppose we have 100 economists, 50 in Aarhus and 50 in Copenhagen
        \end{itemize}
        \item Now suppose that the two universities are hit by demand shocks
        \begin{itemize}
            \item Enrollment is high at Aarhus university sometimes, need more economists
            \item Enrollment is low at Aarhus university sometimes, need fewer economists
        \end{itemize}
        \item If the two universities were in the same location, sometimes when Aarhus university needs fewer economists, Copenhagen university would need more 
        \begin{itemize}
            \item Better for both universities, who need labor, and economists, who need jobs
            \item (or is it?)
        \end{itemize}
    \end{itemize}

\end{frame}

\begin{frame}{Knowledge Spillovers}

    \begin{itemize}
        \item Close to my heart
        \begin{itemize}
            \item Academics spread citations when they move across departments
            \item Can't resist showing you a picture from my paper: 

            http://veryshuai.no-ip.org/downloads/white\_only\_US\_sm.gif
        \end{itemize}
        \item People spread ideas informally and face-to-face
        \begin{itemize}
            \item The business cliché that water-cooler discussions are more important than formal meetings
            \item Recently Yahoo called all telecommuters back
        \end{itemize}
        \item Why are industries clustered within cities?
        \begin{itemize}
            \item Labor market pooling and suppliers can't really explain it
            \item Vernon Hendersen finds that Ad agencies need to be within 300 meters to gain from external economies
        \end{itemize}
    \end{itemize}

\end{frame}

\begin{frame}{Pause}

    \begin{itemize}
        \item We have explained why expect increasing returns at industry-location level
        \item Next
        \begin{itemize}
            \item Setup
            \item Supply, demand, and equilibrium 
            \item Autarchy and trade  
        \end{itemize}
    \end{itemize}

\end{frame}

\begin{frame}

    \begin{itemize}
        \item To start, simplest possible model
        \begin{itemize}
            \item One country
            \item One good
            \item One factor: Labor
            \item \emph{External} increasing returns to scale
            \item That is: industry/location scale economies
        \end{itemize}
    \end{itemize}

\end{frame}

\begin{frame}{The firm's problem}

    \begin{itemize}
        \item Firms are really tiny
        \item They think they do not affect industry wages
        \item In equilibrium, zero-profits again 
        \item What is the equilibrium wage?
    \end{itemize}
    
\end{frame}

\begin{frame}{Equilibrium supply and demand}

    \begin{itemize}
        \item We have been equating relative supply and relative demand
        \item Nothing relative here, we only have one good
        \item Now we are going to equate average cost and demand
        \item Why does avg. cost have to equal price in equilibrium?
        \item Hint: Labor gets everything in this model
    \end{itemize}
    
\end{frame} 
\begin{frame}{Equilibrium supply and demand}

    \begin{itemize}
        \item Last piece: Average cost is falling due to external economies of scale
    \end{itemize}
    \includegraphics[scale=0.20]{ext_equil.png}
    \begin{itemize}
        \item Where does demand curve come from?
    \end{itemize}

\end{frame}

\begin{frame}{Pause}

    \begin{itemize}
        \item Simple 1 x 1 x 1 model
        \item Now add second country
        \item How does trade affect price?
    \end{itemize}

\end{frame}

\begin{frame}{Trade and External Economies}

    \begin{itemize}
        \item Two countries: China and US
        \item Two goods: Buttons and Not-Buttons
        \item One factor: Labor
        \item External economies of scale
    \end{itemize}

\end{frame}

\begin{frame}{No button trade}

    \begin{itemize}
        \item Each country makes its own buttons (and not-buttons)
    \end{itemize}
    \includegraphics[scale=0.20]{buttons_aut.png}

\end{frame}

\begin{frame}{Allow button trade}

    \begin{itemize}
        \item Opening to trade
        \begin{itemize}
            \item China can undercut all American button makers 
            \item True even if very little Labor in American button making
            \item Increasing marginal product of labor (unlike other models)
            \item Implication: China makes all the buttons
        \end{itemize} 
        \item Scale economies
        \begin{itemize}
            \item Cost of producing buttons in China goes down
            \item Cost of producing buttons in US goes up 
            \item Cheaper buttons everywhere!
            \item What about real wages of Chinese workers?
            \item What about real wages of American workers?
        \end{itemize}
    \end{itemize}

\end{frame}

\begin{frame}{Allow button trade}

    \includegraphics[scale=0.20]{button_trade.png}

\end{frame}

\begin{frame}{Implications of external economies model}

    \begin{itemize}
        \item We are going to look at three results
            \begin{enumerate}
                \item The world poverty trap
                \item Losses from trade
                \item Beneficial protectionism
            \end{enumerate}
    \end{itemize}
\end{frame}

\begin{frame}{Poverty traps and models of multiple equilibria}

    \begin{itemize}
        \item Why are some countries poor and others rich?
        \item Destiny or luck?
        \item The classic poverty trap
        \item Models of external economies create a world poverty trap
        \item An industry might be in the wrong country!
    \end{itemize}

\end{frame}

\begin{frame}{Poverty traps and models of multiple equilibria}

    \begin{itemize}
        \item The idea:
        \begin{itemize}
            \item For some reason, only one country produces a good
            \item Scale lowers the price of the good significantly
            \item Another country would be more efficient only if it had all production
            \item At current world price, no firm in other country will enter
        \end{itemize}
    \end{itemize}

\end{frame}

\begin{frame}

    \includegraphics[scale=0.25]{itsatrap.png}

\end{frame}

\begin{frame}{Multiple equilibria}

    \begin{itemize}
        \item Why is an industry in a certain place?
        \begin{itemize}
            \item Often historical accident
        \end{itemize}
        \item Scope for government intervention!
    \end{itemize}

\end{frame}

\begin{frame}{Losses from trade}

    \begin{itemize}
        \item External scale economies can also create losses from trade
        \item Intuition: Same story as wrong equilibrium
        \begin{itemize}
            \item Thailand can make watches cheaper, but Switzerland started first
            \item Now, even if Thailand only produced Thai watches, it would get lower price
            \item Again, low world price keeps firms from entering
        \end{itemize}
    \end{itemize}

\end{frame}

\begin{frame}

    \begin{itemize}
        \item What if Thai government shut down trade?
        \item Who benefits if Thai government shut down trade and then opened it again?
    \end{itemize}
    \includegraphics[scale=0.20]{lossesfromtrade.png}

\end{frame}

\begin{frame}{Learning by doing}

    \begin{itemize}
        \item Suppose that the more a country produces, the lower its cost 
        \item Seems reasonable, lots of studies on WWII America show this
        \item Can give rise to dynamic losses from trade as well
        \begin{itemize}
            \item Two countries, Two goods: Food and Clothing
            \item Rich country has comp. adv. in Food
            \item All learning in agriculture finished
            \item In autarchy, both countries grow 
            \item Suppose Poor specializes in Food after opening to trade
            \item Poor stops growing
        \end{itemize}
        \item Trading dynamic losses for static gains
    \end{itemize}

\end{frame}

\begin{frame}{Learning by doing}

    \begin{itemize}
        \item If there is learning by doing, countries might protect industries until ready for trade
        \item One might call this the Singapore model
        \item Tough part is that it is hard to pick the right industry ex-ante: American solar subsidies
    \end{itemize}

\end{frame}

\begin{frame}{Summary: Chapter 7}

    \begin{itemize}
        \item Implications of external economies model are quite different
        \item \dots and somewhat disturbing
        \item Potential losses from trade
        \item Large beneficial role for government
        \item What is causing these results?
        \begin{itemize}
            \item For one thing, partial equilibrium
        \end{itemize}
    \end{itemize}

\end{frame}

\begin{frame}{Next up: Chapter 8}

    \begin{itemize}
        \item Now, internal increasing returns to scale
        \item The bigger the firm, the lower its cost
        \begin{itemize}
            \item Issue: How many firms if we have two goods? 
            \item Solution: Each firm makes a different product
            \item Implication: Each firm has a monopoly on its product 
        \end{itemize}
        \item Recently very popular, generate firm size distributions
        \item Past chapters: how does trade affect workers and factor owners
        \begin{itemize}
            \item Firms were basically ignored
            \item Side note: If we have L workers, how many free-entry, CRS firms are there?
        \end{itemize}
        \item Main point of chapter: how does trade affect firms?
    \end{itemize}

\end{frame}

\begin{frame}{Monopoly}

    \begin{itemize}
        \item Each firm has a monopoly over its own product
        \item Deconfusifying the textbook review of monopoly, go!
        \begin{itemize}
            \item Price takers vs price makers
            \item Gap between MR and Demand curves
            \item Suppose linear demand: $Q = A - B P$
            \item Average cost vs marginal cost
        \end{itemize}
    \end{itemize}

\end{frame}

\begin{frame}

    \includegraphics[scale=0.25]{monopreview.png}

\end{frame}

\begin{frame}{Monopolistic Competition}

    \begin{itemize}
        \item Monopoly models are easy to analyze
        \item One firm doesn't make an interesting distribution
        \item Oligopoly models are hard
        \begin{itemize}
            \item Pricing decisions depend on beliefs about other firm behavior
            \item With strange beliefs, strange equilibria 
        \end{itemize}
        \item Solution: Monopolistic competition
        \begin{itemize}
            \item Firms are very small
            \item Firms care about market-wide average prices
            \item Individual firm deviations don't change aggregate
        \end{itemize}
    \end{itemize}
            
\end{frame}

\begin{frame}{Active trade literature}

    \begin{itemize}
        \item Starting with Melitz (2002), much active trade literature models of Monop. comp. 
        \item Goal: how does trade affect firms
        \item Monopolistic competition models extremely popular and useful
        \item Differing opinion:
 \emph{The theory of monopolistic competition offers no tools for the analysis of an industry and so no stopping place between the firm at one extreme and general equilibrium at the other. It is therefore incompetent to contribute to the analysis of a host of important problems: the one extreme is too narrow to be of great interest; the other, too broad to permit meaningful generalisations.} -- Milton Friedman (1953), Nobel prize 1976
    \end{itemize}

\end{frame}

\begin{frame}{What we are after}

    \begin{itemize}
        \item How does opening up to trade affect firms?
        \begin{itemize}
            \item Do less productive firms close down?  
            \item Do more firms start operating?
            \item Does each firm get bigger?
        \end{itemize}
        \item How to get there
        \begin{enumerate}
            \item First Monop. Comp. equilibrium in a closed economy
            \item Trade equilibrium with identical firms
            \item Trade equilibrium with heterogenous firms
        \end{enumerate}
    \end{itemize}

\end{frame}

\begin{frame}{Monop. Comp. equilibrium} 

    \begin{itemize}
        \item Plan:
        \begin{itemize}
            \item Show average cost increasing with number of firms
            \item Show price decreasing with number of firms 
            \item Show equilibrium number of firms is intersection (zero profits)
        \end{itemize}
    \end{itemize}

\end{frame}

\begin{frame}{Avg. cost increasing in firm number} 
    \begin{itemize}
        \item Suppose firms face the demand:
            \begin{equation*}
                Q = S \left(\frac{1}{n} - b (P - \overline{P})\right)
            \end{equation*}
        \item If firms are all the same, charge $\overline{P}$
        \item Production of a firm is $Q = \frac{S}{n}$
        \item What happens to production if $n$ increases? 
        \item What happens to the average cost of each firm?
    \end{itemize}

\end{frame}

\begin{frame}{Equil. price decreasing with number of firms} 

    \begin{itemize}
        \item Suppose firms face the demand:
            \begin{equation*}
                Q = S \left(\frac{1}{n} - b (P - \overline{P})\right)
            \end{equation*}
        \item Get price as a function of quantity
        \item Write down firm problem
        \item Write first order condition 
        \item Substitute $Q = \frac{S}{n}$
    \end{itemize}

\end{frame}

\begin{frame}{Price equals average cost} 

    \begin{itemize}
        \item Suppose price were higher than average cost 
            \begin{itemize}
                \item A new firm would enter
            \end{itemize}
        \item Suppose price less than average cost
            \begin{itemize}
                \item Some firm would exit
            \end{itemize}
        \item Must be that price equals average cost
    \end{itemize}

\end{frame}

\begin{frame}{Autarchy equilibrium number of firms}
    
    \includegraphics[scale=0.25]{autequil.png}

\end{frame}

\begin{frame}{Pause}

\begin{itemize}
    \item We have derived the autarchy equilibrium price level and number of firms
    \item Next trade with homogenous firms
\end{itemize}

\end{frame}

\begin{frame}{Opening up to trade}

    \begin{itemize}
        \item First let's show that an increase in market size
        \begin{itemize}
            \item Increases number of active firms (products)
            \item Lowers equilibrium price of each product
        \end{itemize}
    \item Firms produce $\frac{S}{n}$
    \item Market size $S$ increases, each firm produces more
    \item Average cost falls
    \item Price firm's charge is $P = c + \frac{1}{b n}$
    \item Unrelated to market size
    \end{itemize}

\end{frame}

\begin{frame}{Increased market size}
    
    \includegraphics[scale=0.25]{marketsize.png}

\end{frame}

\begin{frame}{Numerical example in book}
    
    \begin{itemize}
        \item Not going to go through it
        \item Point: Opening up to trade 
        \begin{itemize}
            \item Increases number of products available to each consumer
            \item Lowers price of each product
        \end{itemize}
        \item Each firm is a product
        \item Products might be different brands of chocolate 
        \item Intra-industry trade, important part of observed trade
    \end{itemize}

\end{frame}
    
\begin{frame}{Pause}

    \begin{itemize}
        \item Trade with homogenous firms
        \begin{itemize}
            \item Increases available products
            \item Lowers prices
        \end{itemize}
        \item Still not very interesting distribution -- all firms identical!
        \item Now: Heterogenous firms
    \end{itemize}

\end{frame}

\begin{frame}{Heterogenous Firms}

    \begin{itemize}
        \item Firms remain the same in every way
        \item \emph{Except} that they have different marginal costs $c$
        \item Assume that firms learn $c$ \emph{after} they have paid fixed costs 
        \item New equilibrium condition: Expected zero profits (net of fixed cost)
    \end{itemize}

\end{frame}

\begin{frame}{Heterogenous firm equilibrium}

    \begin{itemize}
        \item Firms still produce at point where marginal revenue equals marginal cost
    \end{itemize}
        \includegraphics[scale=0.20]{hetequil.png}

\end{frame}

\begin{frame}{Heterogenous firm equilibrium}

        \begin{itemize}
            \item What is the gap between D and MR?
            \item Who has bigger markup $P - c$?
            \item Do firms with higher marginal cost make more or less profit?
        \end{itemize}
        \includegraphics[scale=0.20]{hetequil_comp.png}

\end{frame}

\begin{frame}{Heterogenous firms and trade}

        \begin{itemize}
            \item Demand:
            \begin{equation*}
                Q = S \left(\frac{1}{n} - b (P - \overline{P})\right)
            \end{equation*}
            \item What is the slope of the demand curve?
            \item What is its intercept?
            \item When a country opens to trade, how does demand change?
        \end{itemize}

\end{frame}
    
\begin{frame}{Heterogenous firms and trade}
    \begin{itemize}
        \item A bit hand-wavy
    \end{itemize}
    \includegraphics[scale=0.20]{hetaftertrade.png}
\end{frame}

\begin{frame}{Heterogenous firms and trade}
    \begin{itemize}
        \item Punchline: Opening up to trade
        \begin{itemize}
            \item Worst firms exit
            \item high-cost firms shrink
            \item low-cost firms grow
        \end{itemize}
        \item Industry more efficient as production is more concentrated in low-cost firms
        \item Effect similar to actual productivity growth
    \end{itemize}
\end{frame}

\begin{frame}{Summary}

    \begin{itemize}
        \item External economies of scale lead to strange results
        \begin{itemize}
            \item Multiple equilibria
            \item Overall losses from trade
            \item Big role for government
        \end{itemize}
        \item Internal economies of scale allow us to consider the distribution of firms
        \begin{itemize}
            \item Opening up to trade has beneficial effects on the firm distribution 
        \end{itemize}
    \end{itemize}

\end{frame}

\begin{frame}

    \begin{itemize}
        \item Coming up
        \begin{itemize}
            \item Wrap up a bit of Chapter 8: trade costs and dumping
            \item Next week the set of chapters on Trade policy (9-12)
            \item Following week we start International Finance
            \item That will mean we did a semester course in 4 weeks!
        \end{itemize}
    \end{itemize}

\end{frame}
\end{document}







